\begin{proposition}[A half-success floor for two-register adaptation]
\label{prop:counter-half-floor}
In the cumulative-counter model, suppose a deterministic two-query
adaptive policy guarantees terminal width at most $\tau$ for every
feasible completion. Then a randomized nonadaptive policy supported on
at most two batches has worst-case precision success at least $1/2$.
Consequently, an adverse-profile upper certificate $U(\tau)=1/2$
establishes the exact optimal randomized-batch success on such an instance.
\end{proposition}
\begin{proof}
If the initial width is already at most $\tau$, a fixed batch trivially
succeeds. Otherwise $E>0$, and let $k$ be the first query of the
guaranteeing policy. Its normalized first answer lies in
$[0,1]$. For each eligible second query $j$, define
\[
 G_j=\{t\in[0,1]:EL_{kj}(t)\leq\tau\}.
\]
Because $L_{kj}$ is affine, $G_j$ is empty, all of $[0,1]$, a closed
prefix interval, or a closed suffix interval. The adaptive guarantee
implies $\bigcup_jG_j=[0,1]$: at each first answer, its chosen second
query must have conditional worst width at most $\tau$.
An execution stopping after the first query can be extended by an unused
second register without increasing width, so such stopping creates no exception.

If some $G_j=[0,1]$, the fixed batch $\{k,j\}$ always succeeds.
Otherwise, choose the longest prefix interval $[0,a]$ and the suffix
interval $[b,1]$ with smallest left endpoint. Both exist because the
union contains the two endpoints, and $a\geq b$ because the union has
no gap. These two intervals therefore cover $[0,1]$.
Select the corresponding fixed batches with probability $1/2$ each.
For any fixed completion, its first answer belongs to at least one of
the two good intervals, and that batch's actual terminal width is bounded
by its conditional worst width. Thus the success probability is at least
$1/2$ at every completion. This matches an upper certificate of $1/2$
when one exists. Endpoint equality counts as success throughout.
\end{proof}

For a deterministic construction, among tied longest prefixes or earliest
suffixes choose the smallest second-query index. This choice depends only
on the coefficient geometry, not on the observed internal registers.
The proposition concerns a fixed two-query budget and this scalar-prefix
information model; it does not assert that arbitrary-budget or arbitrary
query families have the same floor.
